% ═══════════════════════════════════════════════════════════════
%  ch04.tex — فصل ۴
%  الگوی تحلیلی چندبعدی و گونه‌شناسی مداخلات
%  نسخه‌ی اصلاح‌شده: رفع مشکلات RTL، رنگ و جهت
% ═══════════════════════════════════════════════════════════════

\chapter{الگوی تحلیلی چندبعدی و گونه‌شناسی مداخلات}
\label{ch:framework}

\begin{tcolorbox}[quotebox, title={}]
	«بدون چارچوب تحلیلی، تاریخ چیزی جز مجموعه‌ای از حکایات
	پراکنده نیست.»
	
	\hfill ---تدا اسکاچپل\src{\lr{Skocpol, 1979, p.\,xiv}}
\end{tcolorbox}


% ────────────────────────────────────────────────────────────────
\section{درآمد: ضرورت چارچوب تحلیلی}
\label{sec:ch04-intro}

فصل‌های پیشین مفاهیم، تنش‌ها و مبانی حقوقی مداخله‌ی خارجی را
بررسی کردند. اکنون نیاز به ابزاری تحلیلی داریم که بتواند
مطالعات موردی متنوع، از حمله‌ی اعراب به ایران ساسانی تا
مداخله‌ی ناتو در لیبی، را بر مبنای واحد و قابل مقایسه‌ای
ارزیابی کند.

این فصل سه هدف دارد:
\begin{enumerate}[label=\textbf{\arabic*.}, nosep, rightmargin=1em]
	\item ارائه‌ی \concept{الگوی ارزیابی چندبعدی} مداخلات
	(توسعه‌ی ماتریس فصل ۱)؛
	\item تدوین \concept{گونه‌شناسی} جامع انواع مداخلات؛
	\item طراحی \concept{مدل تحلیل پویا} که تعامل بازیگران
	و تحول نتایج در زمان را ترسیم کند.
\end{enumerate}


% ────────────────────────────────────────────────────────────────
\section{الگوی ارزیابی چندبعدی-چندزمانی}
\label{sec:ch04-model}

\subsection{شش بعد ارزیابی}
\label{subsec:ch04-dimensions}

الگوی پیشنهادی بر شش بعد استوار است که هر یک ناظر بر جنبه‌ای
متفاوت از نتایج مداخله‌اند:

\begin{tcolorbox}[mybox,
	title={شش بعد ارزیابی مداخلات خارجی}]
	
	\needspace{4\baselineskip}
	\subsubsection{بعد نظامی-امنیتی}
	\label{subsubsec:ch04-military}
	آیا مداخله به پایان درگیری مسلحانه انجامید؟ آیا ثبات امنیتی
	پایدار ایجاد شد یا خشونت بازتولید شد؟ شاخص‌ها: تلفات نظامی
	و غیرنظامی، مدت درگیری، خلع سلاح، بازسازی نیروهای
	امنیتی\src{\lr{Fortna, 2008, pp.\,1--28}}.
	
	\needspace{4\baselineskip}
	\subsubsection{بعد سیاسی-نهادی}
	\label{subsubsec:ch04-political}
	آیا نهادهای سیاسی مشروع و کارآمد ایجاد شدند؟ آیا حاکمیت ملی
	حفظ شد یا وابستگی سیاسی جایگزین استقلال شد؟ شاخص‌ها: نوع رژیم،
	قانون اساسی، انتخابات، استقلال
	قضایی\src{\lr{Paris, 2004, pp.\,40--55}}.
	
	\needspace{4\baselineskip}
	\subsubsection{بعد اقتصادی}
	\label{subsubsec:ch04-economic}
	آیا مداخله به رشد اقتصادی، بازسازی زیرساخت‌ها و استقلال اقتصادی
	انجامید یا به استثمار منابع و وابستگی اقتصادی؟ شاخص‌ها:
	تولید ناخالص داخلی، نرخ فقر، مالکیت منابع، بدهی
	خارجی\src{\lr{Collier et al., 2003, pp.\,53--80}}.
	
	\needspace{4\baselineskip}
	\subsubsection{بعد حقوق بشری}
	\label{subsubsec:ch04-hr}
	آیا نقض‌های حقوق بشری (نسل‌کشی، شکنجه، تبعیض) متوقف شد؟ آیا
	نهادهای حقوق بشری تأسیس و تقویت شدند؟ شاخص‌ها: گزارش‌های
	سازمان‌های حقوق بشری، وضعیت زندانیان سیاسی، آزادی
	بیان\src{\lr{Donnelly, 2013, pp.\,30--51}}.
	
	\needspace{4\baselineskip}
	\subsubsection{بعد فرهنگی-هویتی}
	\label{subsubsec:ch04-cultural}
	آیا هویت فرهنگی جامعه‌ی هدف حفظ شد؟ آیا تعامل فرهنگی سازنده‌ای
	شکل گرفت یا فرهنگ مداخله‌گر تحمیل شد؟ شاخص‌ها: زبان رسمی،
	آموزش، رسانه، تنوع مذهبی.
	
	\needspace{4\baselineskip}
	\subsubsection{بعد روان‌شناسی اجتماعی}
	\label{subsubsec:ch04-psycho}
	آیا مداخله به حس رهایی و امنیت روانی انجامید یا به ضربه‌ی
	جمعی%
	\footnote{ضربه‌ی جمعی؛ به انگلیسی: \lr{Collective Trauma}}%
	و بی‌اعتمادی مزمن؟ شاخص‌ها: سرمایه‌ی اجتماعی، اعتماد نهادی،
	اختلالات روانی
	جمعی\src{\lr{Volkan, 2004, pp.\,34--78}}.
\end{tcolorbox}

\subsection{سه بازه‌ی زمانی ارزیابی}
\label{subsec:ch04-timeframes}

\begin{tcolorbox}[defbox,
	title={سه بازه‌ی زمانی ارزیابی}]
	\begin{enumerate}[label=\textbf{\arabic*.}, nosep, rightmargin=1em]
		\item \concept{کوتاه‌مدت (۰--۵ سال)}: نتایج فوری مداخله؛
		آیا بحران متوقف شد؟ آیا ثبات اولیه ایجاد شد؟
		\item \concept{میان‌مدت (۵--۲۵ سال)}: نتایج ساختاری؛ آیا
		نهادها تثبیت شدند؟ آیا رشد اقتصادی آغاز شد؟
		\item \concept{بلندمدت (بیش از ۲۵ سال)}: نتایج تمدنی؛ آیا
		جامعه به ثبات پایدار رسید؟ آیا استقلال واقعی حاصل شد؟
	\end{enumerate}
\end{tcolorbox}

\subsection{ماتریس ارزیابی تفصیلی}
\label{subsec:ch04-full-matrix}

ترکیب شش بعد و سه بازه‌ی زمانی، ماتریسی ۱۸ خانه‌ای ایجاد
می‌کند. برای هر مطالعه‌ی موردی، هر خانه با مقیاس پنج‌درجه‌ای
ارزیابی می‌شود:

\begin{table}[htbp]
	\centering
	\caption{ماتریس ارزیابی چندبعدی-چندزمانی (نمونه: ژاپن ۱۹۴۵)}
	\label{tab:ch04-japan-matrix}
	\begin{adjustbox}{max width=\textwidth}
		\begin{tabular}{>{\bfseries\color{PrimaryMid}}p{3cm}
				>{\centering\arraybackslash}p{3.5cm}
				>{\centering\arraybackslash}p{3.5cm}
				>{\centering\arraybackslash}p{3.5cm}}
			\toprule
			\textbf{بُعد} &
			\textbf{کوتاه‌مدت} &
			\textbf{میان‌مدت} &
			\textbf{بلندمدت} \\
			\midrule
			نظامی-امنیتی &
			\cellcolor{green!20}{++ پایان جنگ} &
			\cellcolor{green!20}{++ ثبات امنیتی} &
			\cellcolor{green!25}{++ صلح پایدار} \\
			\midrule
			سیاسی-نهادی &
			\cellcolor{green!15}{+ قانون اساسی} &
			\cellcolor{green!20}{++ دموکراسی} &
			\cellcolor{green!25}{++ تثبیت کامل} \\
			\midrule
			اقتصادی &
			\cellcolor{yellow!20}{$-$ ویرانی} &
			\cellcolor{green!20}{++ معجزه‌ی اقتصادی} &
			\cellcolor{green!25}{++ قدرت سوم} \\
			\midrule
			حقوق بشری &
			\cellcolor{green!15}{+ پایان سرکوب} &
			\cellcolor{green!15}{+ نهادسازی} &
			\cellcolor{green!20}{+ نسبتاً پایدار} \\
			\midrule
			فرهنگی-هویتی &
			\cellcolor{yellow!20}{$\pm$ فشار فرهنگی} &
			\cellcolor{yellow!15}{$\pm$ تنش هویتی} &
			\cellcolor{green!15}{+ ترکیب بومی} \\
			\midrule
			روان‌شناختی &
			\cellcolor{yellow!20}{$\pm$ ضربه و رهایی} &
			\cellcolor{green!15}{+ بازسازی اعتماد} &
			\cellcolor{green!15}{+ نسبتاً مثبت} \\
			\bottomrule
		\end{tabular}
	\end{adjustbox}
\end{table}


% ────────────────────────────────────────────────────────────────
\section{گونه‌شناسی جامع مداخلات}
\label{sec:ch04-typology}

\subsection{پنج معیار طبقه‌بندی}
\label{subsec:ch04-criteria}

مداخلات خارجی را بر اساس پنج معیار می‌توان طبقه‌بندی کرد:

\begin{enumerate}[label=\textbf{معیار \arabic*:}, nosep, rightmargin=1em]
	\item \concept{منشأ تصمیم}: دعوت‌شده / تحمیلی / ترکیبی؛
	\item \concept{ماهیت مداخله‌گر}: دولت / سازمان بین‌المللی /
	بازیگر غیردولتی؛
	\item \concept{نوع هم‌چسبندگی}: ارگانیک / ابزاری / امپریال؛
	\item \concept{شدت مداخله}: فرهنگی / دیپلماتیک / اقتصادی /
	اطلاعاتی / نظامی؛
	\item \concept{هدف اعلامی}: نجات‌بخشی / دفاع / تغییر رژیم /
	گسترش ایدئولوژی / کسب منابع.
\end{enumerate}

\subsection{درخت تصمیم گونه‌شناختی}
\label{subsec:ch04-tree}

\begin{figure}[htbp]
	\centering
	\begin{tikzpicture}[
		>=Stealth,
		every node/.style={font=\small},
		level 1/.style={sibling distance=7cm, level distance=2.5cm},
		level 2/.style={sibling distance=3.5cm, level distance=2.5cm},
		rootnode/.style={
			draw=PrimaryDeep, fill=BgCool, rounded corners=4pt,
			text width=3.2cm, align=center, font=\small\bfseries,
			minimum height=1cm
		},
		leafnode/.style={
			draw=PrimaryMid, fill=BgWarm!50, rounded corners=3pt,
			text width=2.5cm, align=center, font=\scriptsize,
			minimum height=0.9cm
		},
		edge from parent/.style={draw, thick, ->, PrimaryMid}
		]
		
		\node[rootnode] {\rl{مداخله‌ی خارجی}}
		child {
			node[rootnode, fill=BgTeal!40] {\rl{دعوت‌شده}}
			child {
				node[leafnode] {\rl{ارگانیک}\\[1pt]{\rl{(هم‌کیش/هم‌تبار)}}}
				edge from parent node[left, font=\scriptsize, text=AccentTeal]
				{\rl{پیوند تاریخی}}
			}
			child {
				node[leafnode] {\rl{ابزاری}\\[1pt]{\rl{(محاسبه‌ی منفعت)}}}
				edge from parent node[right, font=\scriptsize, text=AccentAmber]
				{\rl{بدون پیوند}}
			}
			edge from parent node[left, font=\scriptsize, text=AccentTeal]
			{\rl{درخواست داخلی}}
		}
		child {
			node[rootnode, fill=red!10] {\rl{تحمیلی}}
			child {
				node[leafnode, fill=red!8] {\rl{بشردوستانه}\\[1pt]{\rl{(\lr{R2P})}}}
				edge from parent node[left, font=\scriptsize, text=PrimaryMid]
				{\rl{توجیه اخلاقی}}
			}
			child {
				node[leafnode, fill=red!15] {\rl{امپریال}\\[1pt]{\rl{(کشورگشایی)}}}
				edge from parent node[right, font=\scriptsize, text=AccentRed]
				{\rl{توسعه‌طلبی}}
			}
			edge from parent node[right, font=\scriptsize, text=AccentRed]
			{\rl{بدون درخواست}}
		};
		
		\node[font=\normalsize\bfseries, text=PrimaryDeep]
		at (0,1.5) {\rl{درخت تصمیم گونه‌شناسی مداخلات}};
	\end{tikzpicture}
	\caption{درخت تصمیم گونه‌شناسی مداخلات خارجی}
	\label{fig:ch04-decision-tree}
\end{figure}

\subsection{جدول گونه‌شناسی تفصیلی}
\label{subsec:ch04-typology-table}

\begin{landscape}
	\begin{table}[htbp]
		\centering
		\caption{گونه‌شناسی تفصیلی مداخلات خارجی با نمونه‌های تاریخی}
		\label{tab:ch04-typology}
		\begin{adjustbox}{max width=\linewidth}
			\begin{tabular}{>{\bfseries\color{PrimaryMid}\small}p{2.3cm}
					>{\small}p{1.8cm} >{\small}p{2cm} >{\small}p{1.8cm}
					>{\small}p{2cm} >{\small}p{2.8cm} >{\small}p{3.2cm}}
				\toprule
				\textbf{گونه} & \textbf{منشأ} & \textbf{هم‌چسبندگی} &
				\textbf{شدت} & \textbf{هدف} & \textbf{نمونه} &
				\textbf{نتیجه} \\
				\midrule
				ارگانیک-رهایی &
				دعوت‌شده &
				هم‌کیشی &
				نظامی &
				رهایی از ستم &
				هند-بنگلادش ۱۹۷۱ &
				موفق: کشور جدید \\
				\midrule
				ارگانیک-الحاقی &
				دعوت‌شده &
				هم‌تباری &
				نظامی &
				وحدت قومی &
				کریمه-روسیه ۲۰۱۴ &
				مناقشه‌آمیز \\
				\midrule
				ابزاری-استراتژیک &
				دعوت‌شده &
				هم‌منفعتی &
				نظامی &
				ثبات/نفوذ &
				سعودی-بحرین ۲۰۱۱ &
				سرکوب اعتراضات \\
				\midrule
				ابزاری-فرصت‌طلبانه &
				دعوت‌شده &
				بدون پیوند &
				نظامی &
				منابع/نفوذ &
				خزعل-بریتانیا &
				ادغام خوزستان \\
				\midrule
				بشردوستانه &
				تحمیلی &
				--- &
				نظامی &
				حمایت غیرنظامیان &
				ناتو-کوزوو ۱۹۹۹ &
				نسبتاً موفق \\
				\midrule
				امپریال-ایدئولوژیک &
				تحمیلی &
				ادعایی &
				نظامی &
				تغییر رژیم &
				شوروی-افغانستان ۱۹۷۹ &
				شکست \\
				\midrule
				میسیونری-فرهنگی &
				نرم &
				هم‌کیشی &
				فرهنگی &
				تغییر ارزش‌ها &
				میسیونرها در آفریقا &
				تغییرات عمیق \\
				\midrule
				صادراتی-ایدئولوژیک &
				نرم تا سخت &
				هم‌ایدئولوژی &
				ترکیبی &
				صدور ایدئولوژی &
				صدور انقلاب &
				مختلط \\
				\midrule
				پان‌ناسیونالیستی &
				ترکیبی &
				هم‌تباری &
				ترکیبی &
				وحدت ملی &
				پان‌ترکیسم &
				عمدتاً ناموفق \\
				\bottomrule
			\end{tabular}
		\end{adjustbox}
	\end{table}
\end{landscape}


% ────────────────────────────────────────────────────────────────
\section{مدل تحلیل پویای مداخله}
\label{sec:ch04-dynamic}

مداخلات خارجی فرایندهای ایستا نیستند. آن‌ها در طول زمان تحول
می‌یابند و بازیگران مختلف بر یکدیگر تأثیر می‌گذارند. برای
مدل‌سازی این پویایی،
\concept{مدل سه‌لایه‌ی بازیگر-ساختار-نتیجه}
را پیشنهاد می‌کنیم:

\subsection{لایه‌ی اول: بازیگران}
\label{subsec:ch04-actors}

\begin{tcolorbox}[mybox,
	title={بازیگران اصلی در فرایند مداخله}]
	\begin{enumerate}[label=\textbf{\arabic*.}, nosep, rightmargin=1em]
		\item \concept{دعوت‌کننده(گان)}: نخبگان سیاسی، اقلیت قومی-مذهبی،
		مرزنشینان، جنبش اجتماعی، یا بخشی از مردم عادی؛
		\item \concept{مداخله‌گر}: دولت خارجی، سازمان بین‌المللی،
		بازیگر غیردولتی، یا ائتلاف؛
		\item \concept{دولت/ساختار مرکزی}: رژیم حاکم بر جامعه‌ی هدف؛
		\item \concept{جمعیت عام}: مردم عادی که اغلب نه دعوت‌کننده‌اند
		و نه تصمیم‌گیرنده، ولی بیشترین تأثیر را می‌پذیرند؛
		\item \concept{ناظران بین‌المللی}: دولت‌ها و سازمان‌هایی که
		مستقیم درگیر نیستند ولی واکنش نشان می‌دهند؛
		\item \concept{بازیگران رقیب}: نیروهای خارجی دیگری که منافع
		متضاد با مداخله‌گر دارند.
	\end{enumerate}
\end{tcolorbox}

\subsection{لایه‌ی دوم: ساختارها و زمینه‌ها}
\label{subsec:ch04-structures}

\begin{table}[htbp]
	\centering
	\caption{ساختارها و زمینه‌های مؤثر بر نتیجه‌ی مداخله}
	\label{tab:ch04-structures}
	\begin{adjustbox}{max width=\textwidth}
		\begin{tabular}{>{\bfseries\color{PrimaryMid}}p{2.8cm}
				p{4.5cm} p{4.5cm}}
			\toprule
			\textbf{ساختار} & \textbf{توصیف} & \textbf{تأثیر بر نتیجه} \\
			\midrule
			ساختار قومی-مذهبی &
			درجه‌ی همگنی یا چندپارگی جامعه &
			چندپارگی بیشتر = ریسک بالاتر شکست \\
			\midrule
			ظرفیت دولتی &
			توانایی بوروکراسی و نهادهای محلی &
			ظرفیت بالاتر = شانس بیشتر نهادسازی \\
			\midrule
			جغرافیای سیاسی &
			موقعیت ژئوپولیتیک و همسایگان &
			همسایگان مداخله‌جو = بی‌ثباتی بیشتر \\
			\midrule
			تاریخ مداخلات &
			تجربه‌ی پیشین جامعه با مداخله‌ی خارجی &
			تجربه‌ی منفی = بی‌اعتمادی و مقاومت \\
			\midrule
			اقتصاد سیاسی &
			منابع طبیعی، وابستگی اقتصادی &
			منابع غنی = انگیزه‌ی بیشتر مداخله‌گر \\
			\midrule
			نظم بین‌المللی &
			ساختار نظام جهانی &
			تک‌قطبی = مداخله‌ی آسان‌تر \\
			\bottomrule
		\end{tabular}
	\end{adjustbox}
\end{table}

\subsection{لایه‌ی سوم: مکانیسم‌ها و نتایج}
\label{subsec:ch04-mechanisms}

\begin{figure}[htbp]
	\centering
	\begin{tikzpicture}[
		>=Stealth,
		every node/.style={font=\small},
		phase/.style={
			draw=PrimaryMid, fill=#1, rounded corners=4pt,
			minimum width=3cm, minimum height=1.6cm,
			text width=2.6cm, align=center, font=\small
		},
		mechanism/.style={
			draw=AccentGold, fill=BgWarm!40, rounded corners=2pt,
			text width=2cm, align=center, font=\scriptsize
		}
		]
		
		% فازها
		\node[phase=BgCool!50] (pre) at (0,0)
		{\textbf{\rl{پیشا-مداخله}}\\[2pt]%
			\rl{بحران، ستم}\\[1pt]%
			\rl{درخواست/تصمیم}};
		
		\node[phase=BgTeal!40] (int) at (5.5,0)
		{\textbf{\rl{فاز مداخله}}\\[2pt]%
			\rl{عملیات نظامی}\\[1pt]%
			\rl{فشار سیاسی/اقتصادی}};
		
		\node[phase=BgWarm!50] (post) at (11,0)
		{\textbf{\rl{پسا-مداخله}}\\[2pt]%
			\rl{بازسازی، نهادسازی}\\[1pt]%
			\rl{یا فروپاشی و بازگشت}};
		
		% فلش‌های اصلی
		\draw[->, thick, PrimaryMid] (pre) -- (int);
		\draw[->, thick, PrimaryMid] (int) -- (post);
		
		% بازخورد
		\draw[->, thick, AccentRed, dashed, bend right=35]
		(post.south) to node[below, font=\scriptsize, text=AccentRed,
		text width=3cm, align=center]
		{\rl{بازگشت بحران}\\[1pt]{\scriptsize\rl{(چرخه‌ی معیوب)}}}
		(pre.south);
		
		% نتایج
		\node[mechanism, fill=green!15] (r1) at (8,-3)
		{\rl{نتیجه‌ی مثبت:}\\[1pt]{\rl{رهایی و ثبات}}};
		\node[mechanism, fill=red!10] (r2) at (14,-3)
		{\rl{نتیجه‌ی منفی:}\\[1pt]{\rl{وابستگی/فروپاشی}}};
		
		\draw[->, thin, AccentGreen] (post.south) -- ++(0,-0.8) -| (r1.north);
		\draw[->, thin, AccentRed] (post.south) -- ++(0,-0.8) -| (r2.north);
		
		% عنوان
		\node[font=\normalsize\bfseries, text=PrimaryDeep]
		at (5.5,2.5) {\rl{مدل سه‌فازی تحلیل پویای مداخله}};
	\end{tikzpicture}
	\caption{مدل سه‌فازی تحلیل پویای مداخله‌ی خارجی}
	\label{fig:ch04-dynamic-model}
\end{figure}


% ────────────────────────────────────────────────────────────────
\section{مدل هم‌چسبندگی: ارگانیک در برابر ابزاری}
\label{sec:ch04-adhesion-model}

\subsection{شاخص‌های تعیین نوع هم‌چسبندگی}
\label{subsec:ch04-adhesion-indicators}

\begin{table}[htbp]
	\centering
	\caption{شاخص‌های تفکیک نوع هم‌چسبندگی}
	\label{tab:ch04-adhesion}
	\begin{adjustbox}{max width=\textwidth}
		\begin{tabular}{>{\bfseries\color{PrimaryMid}}p{2.8cm}
				p{3.2cm} p{3.2cm} p{3.2cm}}
			\toprule
			\textbf{شاخص} & \textbf{ارگانیک} &
			\textbf{ابزاری} & \textbf{امپریال} \\
			\midrule
			سابقه‌ی تاریخی &
			چندین سده یا بیشتر &
			کوتاه‌مدت یا ناموجود &
			ممکن است ادعایی باشد \\
			\midrule
			عمق فرهنگی-زبانی &
			زبان/مذهب/آداب مشترک &
			تفاوت‌های اساسی &
			--- \\
			\midrule
			جهت‌گیری پس از مداخله &
			ادغام یا خودمختاری &
			حفظ فاصله &
			تصرف و سلطه \\
			\midrule
			توزیع منافع &
			نسبتاً متقارن &
			به نفع یک طرف &
			به‌شدت نامتقارن \\
			\midrule
			مشارکت محلی &
			بالا &
			متوسط تا پایین &
			حداقلی یا نمایشی \\
			\midrule
			برنامه‌ی خروج &
			ادغام یا خروج واقعی &
			مبهم &
			بدون خروج \\
			\bottomrule
		\end{tabular}
	\end{adjustbox}
\end{table}


% ────────────────────────────────────────────────────────────────
\section{الگوی پیش‌زمینه‌سازی: از میسیونری تا نظامی}
\label{sec:ch04-precursor-model}

\subsection{زنجیره‌ی تصاعد}
\label{subsec:ch04-escalation}

\begin{figure}[htbp]
	\centering
	\begin{tikzpicture}[
		>=Stealth,
		stage/.style={
			draw=PrimaryMid, fill=#1, rounded corners=3pt,
			minimum width=2.2cm, minimum height=1.4cm,
			text width=2cm, align=center, font=\small
		}
		]
		
		\node[stage=BgTeal!30] (s1) at (0,0)
		{\textbf{\rl{فرهنگی}}\\[1pt]%
			{\scriptsize\rl{میسیونری}}};
		
		\node[stage=BgWarm!40] (s2) at (3.2,0)
		{\textbf{\rl{سیاسی}}\\[1pt]%
			{\scriptsize\rl{دیپلماسی}}};
		
		\node[stage=BgCool!50] (s3) at (6.4,0)
		{\textbf{\rl{اقتصادی}}\\[1pt]%
			{\scriptsize\rl{تحریم/کمک}}};
		
		\node[stage=yellow!20] (s4) at (9.6,0)
		{\textbf{\rl{اطلاعاتی}}\\[1pt]%
			{\scriptsize\rl{سایبری/کودتا}}};
		
		\node[stage=red!12] (s5) at (12.8,0)
		{\textbf{\rl{نظامی}}\\[1pt]%
			{\scriptsize\rl{حمله/اشغال}}};
		
		% فلش‌های تصاعد
		\draw[->, very thick, AccentGold] (s1) -- (s2);
		\draw[->, very thick, AccentGold] (s2) -- (s3);
		\draw[->, very thick, AccentGold] (s3) -- (s4);
		\draw[->, very thick, AccentGold] (s4) -- (s5);
		
		% فلش شدت
		\draw[->, very thick, AccentRed] (0,-1.8) -- (12.8,-1.8);
		\node[font=\small, text=AccentRed] at (6.4,-2.2)
		{\rl{افزایش شدت و هزینه}};
		
		% فلش برگشت‌پذیری
		\draw[<-, very thick, AccentGreen] (0,-2.8) -- (12.8,-2.8);
		\node[font=\small, text=AccentGreen] at (6.4,-3.2)
		{\rl{کاهش برگشت‌پذیری}};
		
		\node[font=\normalsize\bfseries, text=PrimaryDeep]
		at (6.4,2) {\rl{زنجیره‌ی تصاعد مداخله‌ی خارجی}};
	\end{tikzpicture}
	\caption{زنجیره‌ی تصاعد: از فعالیت فرهنگی تا مداخله‌ی نظامی}
	\label{fig:ch04-escalation}
\end{figure}

\subsection{سه الگوی رابطه‌ی فعالیت فرهنگی و مداخله‌ی سخت}
\label{subsec:ch04-missionary}

\begin{tcolorbox}[comparebox,
	title={سه الگوی رابطه‌ی فعالیت فرهنگی و مداخله‌ی سخت}]
	\begin{enumerate}[label=\textbf{\arabic*.}, nosep, rightmargin=1em]
		\item \concept{الگوی پیشقراول}%
		\footnote{الگوی پیشقراول؛ به انگلیسی: \lr{Vanguard Model}}:
		فعالیت فرهنگی آگاهانه مقدمه‌ی استعمار است.
		نمونه: میسیونرهای مسیحی در آفریقا پیش از تقسیم‌بندی
		برلین\src{\lr{Comaroff \& Comaroff, 1991}}؛
		\item \concept{الگوی موازی}%
		\footnote{الگوی موازی؛ به انگلیسی: \lr{Parallel Model}}:
		فعالیت فرهنگی و مداخله‌ی سخت مستقل از هم ولی همزمان.
		نمونه: ترویج لیبرال-دموکراسی آمریکایی و عملیات سیا؛
% ═══════════════════════════════════════════════════════════════
%  ch04.tex — ادامه (از الگوی مستقل میسیونری)
% ═══════════════════════════════════════════════════════════════

\item \concept{الگوی مستقل}%
\footnote{الگوی مستقل؛ به انگلیسی: \lr{Independent Model}}:
فعالیت فرهنگی بدون بعد استعماری و حتی در تعارض با قدرت‌های
استعماری. نمونه: فعالیت‌های روشنفکران مشروطه‌خواه ایرانی
در تفلیس و باکو\src{\lr{Adamiyat, 1961, pp.\,34--56}}.
\end{enumerate}
\end{tcolorbox}

\begin{tcolorbox}[wavebox,
title={نکته‌ی تحلیلی: تفکیک میسیونری از استعمار}]
تفکیک فعالیت‌های فرهنگی-ایدئولوژیک از پروژه‌ی استعماری نیازمند
بررسی سه عامل است:
\begin{enumerate}[label=\textbf{\alph*.}, nosep, rightmargin=1em]
\item \concept{نیت}: آیا فعالیت فرهنگی آگاهانه در خدمت
پروژه‌ی سلطه قرار دارد یا مستقل از آن است؟
\item \concept{ساختار}: آیا فعالیت فرهنگی از حمایت مالی و
لجستیکی دولت استعمارگر بهره‌مند است؟
\item \concept{نتیجه}: آیا فعالیت فرهنگی عملاً به تسهیل سلطه‌ی
سیاسی-اقتصادی منجر شده؟
\end{enumerate}
تنها هنگامی که هر سه عامل مثبت باشند، می‌توان فعالیت فرهنگی را
ذیل مقوله‌ی استعمار طبقه‌بندی کرد. در سایر موارد، باید آن‌ها را
مقوله‌ی مستقلی دانست\src{\lr{Stanley, 1990, pp.\,45--72}}.
\end{tcolorbox}


% ────────────────────────────────────────────────────────────────
\section{الگوی ارزیابی ژئوپولیتیک مداخله}
\label{sec:ch04-geopolitical}

\subsection{نتایج ژئوپولیتیک: تشکیل کشور، تقسیم مرزها، یا ادغام}
\label{subsec:ch04-geopolitical-outcomes}

مداخلات خارجی ممکن است به یکی از سه نتیجه‌ی ژئوپولیتیک بینجامند:

\begin{table}[htbp]
\centering
\caption{نتایج ژئوپولیتیک مداخلات خارجی}
\label{tab:ch04-geopolitical}
\begin{adjustbox}{max width=\textwidth}
\begin{tabular}{>{\bfseries\color{PrimaryMid}}p{2.5cm}
		p{3.2cm} p{3.2cm} p{3.5cm}}
	\toprule
	\textbf{نتیجه} & \textbf{شرایط} & \textbf{نمونه‌ی موفق} &
	\textbf{نمونه‌ی ناموفق} \\
	\midrule
	تشکیل کشور جدید &
	حمایت بین‌المللی + انسجام داخلی &
	بنگلادش ۱۹۷۱ &
	سودان جنوبی ۲۰۱۱ \\
	\midrule
	تقسیم و بازترسیم مرزها &
	ضعف دولت مرکزی + رقابت قدرت‌ها &
	تجزیه‌ی یوگسلاوی &
	تجزیه‌ی عثمانی (مختلط) \\
	\midrule
	ادغام و حفظ تمامیت &
	دولت مرکزی قوی + حمایت بین‌المللی &
	نیجریه پس از بیافرا &
	سوریه (حفظ صوری) \\
	\bottomrule
\end{tabular}
\end{adjustbox}
\end{table}


% ────────────────────────────────────────────────────────────────
\section{نقش روان‌شناسی اجتماعی در الگوی تحلیلی}
\label{sec:ch04-psychology}

\subsection{روان‌شناسی نجات‌طلبی}
\label{subsec:ch04-rescue-psychology}

\thinker{وامیک ولکان}%
\footnote{%
وامیک ولکان (\lr{Vamık D.\,Volkan, 1932–})؛%
روان‌پزشک و نظریه‌پرداز روابط بین‌گروهی ترکیه‌ای-آمریکایی.%
}%
مفهوم \concept{آسیب‌های انتخابی}%
\footnote{آسیب‌های انتخابی؛ به انگلیسی: \lr{Chosen Traumas}}%
را مطرح کرده: رویدادهای دردناک تاریخی که هویت جمعی یک گروه حول
آن‌ها شکل می‌گیرد و از نسلی به نسل دیگر منتقل
می‌شود\src{\lr{Volkan, 2004, pp.\,47--78}}.

این آسیب‌های انتخابی نقش مهمی در پویایی نجات‌طلبی دارند:
گروهی که خاطره‌ی جمعی آن حول رنج و ستم شکل گرفته، مستعد
جستجوی ناجی بیرونی است؛ حتی اگر تجربه‌ی تاریخی نشان دهد
که ناجیان بیرونی اغلب خود به ستمگران جدید تبدیل شده‌اند.

\subsection{ایده‌آل‌سازی ناجی خارجی}
\label{subsec:ch04-idealization}

\begin{tcolorbox}[defbox,
title={مفهوم: ایده‌آل‌سازی ناجی خارجی}]
در روان‌شناسی اجتماعی، پدیده‌ای وجود دارد که می‌توان آن را
\concept{ایده‌آل‌سازی ناجی خارجی}%
\footnote{ایده‌آل‌سازی ناجی خارجی؛ به انگلیسی:
\lr{Idealization of the External Savior}}%
نامید: گروه تحت ستم، از شدت ناامیدی از وضع داخلی، بیگانه
را ایده‌آل می‌کند و نیات و توانایی‌های او را بیش از واقع تصور
می‌کند. این پدیده:
\begin{itemize}[nosep, rightmargin=1em]
\item با \concept{عقلانیت محدود} سازگار است: در بحران،
اطلاعات ناقص و هیجانات بر تصمیم‌گیری غلبه
دارند\src{\lr{Simon, 1957}}؛
\item با \concept{ناهماهنگی شناختی}%
\footnote{ناهماهنگی شناختی؛ به انگلیسی:
	\lr{Cognitive Dissonance}}
قابل تبیین است: برای توجیه تصمیم به طلب کمک خارجی، مزایای
آن بزرگ و معایبش کوچک جلوه داده
می‌شوند\src{\lr{Festinger, 1957, pp.\,1--31}}؛
\item با \concept{نظریه‌ی هویت اجتماعی}%
\footnote{نظریه‌ی هویت اجتماعی؛ به انگلیسی:
	\lr{Social Identity Theory}}
مرتبط است: «خودی» بد و «بیگانه» خوب تلقی
می‌شود\src{\lr{Tajfel \& Turner, 1979, pp.\,33--47}}.
\end{itemize}
\end{tcolorbox}

\subsection{چرخه‌ی روان‌شناختی مداخله}
\label{subsec:ch04-psych-cycle}

\begin{figure}[htbp]
\centering
\begin{tikzpicture}[
>=Stealth,
every node/.style={font=\small},
stage/.style={
	draw=PrimaryMid, fill=#1, rounded corners=4pt,
	text width=2.4cm, align=center, minimum height=1.2cm
}
]

\node[stage=BgCool!40] (s1) at (90:3.8)
{\textbf{\rl{ناامیدی}}\\[1pt]%
	{\scriptsize\rl{از وضع داخلی}}};

\node[stage=BgTeal!40] (s2) at (18:3.8)
{\textbf{\rl{ایده‌آل‌سازی}}\\[1pt]%
	{\scriptsize\rl{ناجی خارجی}}};

\node[stage=BgWarm!40] (s3) at (306:3.8)
{\textbf{\rl{مداخله و}}\\[1pt]%
	{\scriptsize\rl{انتظارات بالا}}};

\node[stage=red!10] (s4) at (234:3.8)
{\textbf{\rl{سرخوردگی}}\\[1pt]%
	{\scriptsize\rl{نتایج ناامیدکننده}}};

\node[stage=yellow!20] (s5) at (162:3.8)
{\textbf{\rl{بی‌اعتمادی و}}\\[1pt]%
	{\scriptsize\rl{خشم مزمن}}};

% فلش‌ها
\draw[->, very thick, AccentGold] (s1) -- (s2);
\draw[->, very thick, AccentGold] (s2) -- (s3);
\draw[->, very thick, AccentRed] (s3) -- (s4);
\draw[->, very thick, AccentRed] (s4) -- (s5);
\draw[->, very thick, NeutralDark, dashed] (s5) -- (s1);

% مرکز
\node[font=\normalsize\bfseries, text=PrimaryDeep,
text width=2.8cm, align=center]
at (0,0) {\rl{چرخه‌ی}\\[2pt]{\rl{روان‌شناختی}}\\[2pt]{\rl{مداخله}}};
\end{tikzpicture}
\caption{چرخه‌ی روان‌شناختی نجات‌طلبی و سرخوردگی}
\label{fig:ch04-psych-cycle}
\end{figure}

\begin{tcolorbox}[wavebox,
title={تبیین چرخه‌ی روان‌شناختی}]
\begin{enumerate}[label=\textbf{گام \arabic*:}, nosep, rightmargin=1em]
\item \concept{ناامیدی}: فشار مداوم رژیم، نبود افق تغییر
داخلی، احساس عجز و درماندگی؛
\item \concept{ایده‌آل‌سازی}: نیروی خارجی به‌مثابه‌ی منجی
ترسیم می‌شود؛ اطلاعات مثبت بزرگ‌نمایی و اطلاعات منفی
کوچک‌نمایی می‌شوند؛
\item \concept{مداخله}: نیروی خارجی وارد می‌شود؛ انتظارات
بالا و ناواقع‌بینانه شکل گرفته؛
\item \concept{سرخوردگی}: نتایج مداخله با انتظارات تطابق
ندارد؛ مداخله‌گر منافع خود را دنبال می‌کند؛ هزینه‌های
انسانی بالاست؛
\item \concept{بی‌اعتمادی}: خشم و ناامیدی نه‌تنها نسبت به
مداخله‌گر بلکه نسبت به کل نهادهای بیرونی شکل می‌گیرد؛
این بی‌اعتمادی به نسل بعد منتقل شده و زمینه‌ی چرخه‌ی
بعدی را فراهم می‌کند.
\end{enumerate}
\end{tcolorbox}

\subsection{شکستن چرخه: شرایط روان‌شناختی موفقیت}
\label{subsec:ch04-break-cycle}

\thinker{هربرت کلمن}%
\footnote{%
هربرت کلمن (\lr{Herbert C.\,Kelman, 1927--2022})؛%
روان‌شناس اجتماعی آمریکایی و پیشگام حل تعارض تعاملی.%
}%
در پژوهش‌هایش درباره‌ی حل تعارض اسرائیلی-فلسطینی نشان داده
که شکستن چرخه‌ی بی‌اعتمادی نیازمند
\concept{بازشناسایی متقابل}%
\footnote{بازشناسایی متقابل؛ به انگلیسی: \lr{Mutual Recognition}}
است: هر طرف باید هویت، رنج و حقوق طرف مقابل را به رسمیت
بشناسد\src{\lr{Kelman, 2005, pp.\,111--129}}.

در بافت مداخله‌ی خارجی، این به معنای آن است که:
\begin{enumerate}[label=\textbf{\arabic*.}, nosep, rightmargin=1em]
\item مداخله‌گر باید رنج‌های ناشی از مداخله را به رسمیت بشناسد؛
\item جامعه‌ی هدف باید بتواند روایت خود را از تجربه‌ی مداخله
بازگو و بازسازی کند؛
\item فرایند عدالت انتقالی%
\footnote{عدالت انتقالی؛ به انگلیسی:
\lr{Transitional Justice}}
باید اجرا شود\src{\lr{Teitel, 2000, pp.\,3--26}}.
\end{enumerate}


% ────────────────────────────────────────────────────────────────
\section{الگوی ترکیبی نهایی: پنج‌ضلعی تحلیل مداخله}
\label{sec:ch04-pentagon}

با ترکیب تمام ابعاد مطرح‌شده، الگوی نهایی تحلیل مداخلات خارجی
را در قالب یک \concept{پنج‌ضلعی تحلیلی} ارائه می‌کنیم:

\begin{figure}[htbp]
\centering
\begin{tikzpicture}[
>=Stealth,
every node/.style={font=\small}
]

% پنج‌ضلعی — صریح (بدون foreach)
\node[draw=AccentTeal, fill=AccentTeal!10, rounded corners=4pt,
minimum width=2.5cm, minimum height=1cm,
align=center, font=\small\bfseries, text=AccentTeal]
(n1) at (90:4.5) {\rl{بازیگران}};

\node[draw=PrimaryMid, fill=PrimaryMid!10, rounded corners=4pt,
minimum width=2.5cm, minimum height=1cm,
align=center, font=\small\bfseries, text=PrimaryMid]
(n2) at (162:4.5) {\rl{ساختارها}};

\node[draw=AccentAmber, fill=AccentAmber!10, rounded corners=4pt,
minimum width=2.5cm, minimum height=1cm,
align=center, font=\small\bfseries, text=AccentAmber]
(n3) at (234:4.5) {\rl{مکانیسم‌ها}};

\node[draw=AccentGreen, fill=AccentGreen!10, rounded corners=4pt,
minimum width=2.5cm, minimum height=1cm,
align=center, font=\small\bfseries, text=AccentGreen]
(n4) at (306:4.5) {\rl{نتایج}};

\node[draw=AccentRed, fill=AccentRed!10, rounded corners=4pt,
minimum width=2.5cm, minimum height=1cm,
align=center, font=\small\bfseries, text=AccentRed]
(n5) at (18:4.5) {\rl{بازخورد}};

% اتصالات اصلی
\draw[thick, NeutralDark] (n1) -- (n2);
\draw[thick, NeutralDark] (n2) -- (n3);
\draw[thick, NeutralDark] (n3) -- (n4);
\draw[thick, NeutralDark] (n4) -- (n5);
\draw[thick, NeutralDark] (n5) -- (n1);

% اتصالات متقاطع
\draw[thin, NeutralMid, dashed] (n1) -- (n3);
\draw[thin, NeutralMid, dashed] (n1) -- (n4);
\draw[thin, NeutralMid, dashed] (n2) -- (n4);
\draw[thin, NeutralMid, dashed] (n2) -- (n5);
\draw[thin, NeutralMid, dashed] (n3) -- (n5);

% مرکز
\node[draw=PrimaryDeep, fill=BgCool, circle,
minimum size=2.2cm, align=center, font=\bfseries,
text=PrimaryDeep]
at (0,0) {\rl{ارزیابی}\\[2pt]{\rl{چندبعدی-}}\\[2pt]{\rl{چندزمانی}}};

% عنوان
\node[font=\large\bfseries, text=PrimaryDeep]
at (0,6.5) {\rl{پنج‌ضلعی تحلیل مداخله‌ی خارجی}};
\end{tikzpicture}
\caption{الگوی ترکیبی نهایی: پنج‌ضلعی تحلیل مداخله}
\label{fig:ch04-pentagon}
\end{figure}

\begin{tcolorbox}[mybox,
title={توضیح پنج ضلع}]
\begin{enumerate}[label=\textbf{\arabic*.}, nosep, rightmargin=1em]
\item \concept{بازیگران}: چه کسانی درگیرند؟ دعوت‌کننده، مداخله‌گر،
دولت مرکزی، جمعیت عام، ناظران، رقبا؛
\item \concept{ساختارها}: بافت قومی-مذهبی، ظرفیت دولتی، جغرافیای
سیاسی، اقتصاد سیاسی، نظم بین‌المللی؛
\item \concept{مکانیسم‌ها}: زنجیره‌ی تصاعد، نوع هم‌چسبندگی،
سه فاز مداخله (پیشا، حین، پسا)؛
\item \concept{نتایج}: ارزیابی شش‌بعدی-سه‌زمانی (ماتریس ۱۸
خانه‌ای)؛ نتایج ژئوپولیتیک (تشکیل کشور، تقسیم مرزها، ادغام)؛
\item \concept{بازخورد}: چرخه‌ی روان‌شناختی؛ اثر سابقه؛ بازتولید
خشونت یا بازسازی اعتماد.
\end{enumerate}

هر مطالعه‌ی موردی در فصل‌های آتی بر اساس این پنج‌ضلعی تحلیل
خواهد شد. این چارچوب امکان مقایسه‌ی تطبیقی منسجم میان موارد
بسیار متنوع، از حمله‌ی اعراب به ایران ساسانی تا مداخله‌ی
ناتو در لیبی، را فراهم می‌آورد.
\end{tcolorbox}


% ────────────────────────────────────────────────────────────────
\section{خلاصه‌ی ابزارهای تحلیلی}
\label{sec:ch04-tools-summary}

\begin{table}[htbp]
\centering
\caption{خلاصه‌ی ابزارهای تحلیلی فصل ۴}
\label{tab:ch04-tools-summary}
\begin{adjustbox}{max width=\textwidth}
\begin{tabular}{>{\bfseries\color{PrimaryMid}}p{3.5cm}
		p{4.5cm} p{4.5cm}}
	\toprule
	\textbf{ابزار} & \textbf{کاربرد} & \textbf{مرجع شکل/جدول} \\
	\midrule
	ماتریس ۱۸ خانه‌ای &
	ارزیابی منسجم شش‌بعدی-سه‌زمانی &
	جدول \ref{tab:ch04-japan-matrix} \\
	\midrule
	گونه‌شناسی ده‌گانه &
	طبقه‌بندی مداخلات بر اساس ۵ معیار &
	جدول \ref{tab:ch04-typology} \\
	\midrule
	درخت تصمیم &
	تفکیک سریع نوع مداخله &
	شکل \ref{fig:ch04-decision-tree} \\
	\midrule
	مدل سه‌فازی &
	تحلیل پویای مداخله در زمان &
	شکل \ref{fig:ch04-dynamic-model} \\
	\midrule
	شاخص‌های هم‌چسبندگی &
	تفکیک ارگانیک/ابزاری/امپریال &
	جدول \ref{tab:ch04-adhesion} \\
	\midrule
	زنجیره‌ی تصاعد &
	تحلیل مسیر از نرم به سخت &
	شکل \ref{fig:ch04-escalation} \\
	\midrule
	سه الگوی فرهنگی &
	رابطه‌ی فعالیت فرهنگی و استعمار &
	بخش \ref{subsec:ch04-missionary} \\
	\midrule
	چرخه‌ی روان‌شناختی &
	تحلیل پویایی نجات‌طلبی-سرخوردگی &
	شکل \ref{fig:ch04-psych-cycle} \\
	\midrule
	پنج‌ضلعی تحلیل &
	چارچوب ترکیبی نهایی &
	شکل \ref{fig:ch04-pentagon} \\
	\bottomrule
\end{tabular}
\end{adjustbox}
\end{table}


% ────────────────────────────────────────────────────────────────
\section{جمع‌بندی فصل و آمادگی برای مطالعات موردی}
\label{sec:ch04-summary}

در این فصل، مجموعه‌ی کاملی از ابزارهای تحلیلی برای بررسی
مطالعات موردی فراهم شد. این ابزارها عبارت‌اند از:

\begin{enumerate}[label=\textbf{\arabic*.}, nosep, rightmargin=1em]
\item \concept{ماتریس ارزیابی ۱۸ خانه‌ای} (شش بعد ضربدر سه
بازه‌ی زمانی) برای ارزیابی منسجم نتایج؛
\item \concept{گونه‌شناسی ده‌گانه} بر اساس پنج معیار
(منشأ، ماهیت مداخله‌گر، هم‌چسبندگی، شدت، هدف)؛
\item \concept{مدل سه‌فازی پویا} (پیشا-مداخله، مداخله،
پسا-مداخله) با امکان بازخورد و چرخه‌ی معیوب؛
\item \concept{مدل هم‌چسبندگی} برای تفکیک مداخلات ارگانیک،
ابزاری و امپریال؛
\item \concept{زنجیره‌ی تصاعد} و سه الگوی رابطه‌ی فعالیت
فرهنگی با مداخله‌ی سخت؛
\item \concept{مدل روان‌شناختی} چرخه‌ی نجات‌طلبی و سرخوردگی
و شرایط شکستن آن؛
\item \concept{پنج‌ضلعی تحلیل} به‌عنوان چارچوب ترکیبی نهایی.
\end{enumerate}

از بخش دوم کتاب (فصل ۵ به بعد)، این ابزارها بر مطالعات موردی
اعمال خواهند شد؛ ابتدا بر تاریخ ایران، از دوره‌ی هخامنشی تا
دوره‌ی معاصر، و سپس بر تجربیات جهانی در اروپا، خاورمیانه،
آفریقا، آسیا و آمریکای لاتین.

\begin{tcolorbox}[mybox,
title={پرسش راهنمای مطالعات موردی}]
برای هر مطالعه‌ی موردی، پنج پرسش اصلی پاسخ داده خواهد شد:

\begin{enumerate}[label=\textbf{پرسش \arabic*:}, nosep, rightmargin=1em]
\item \textbf{بازیگران}: چه کسی دعوت کرد، چه کسی مداخله کرد
و چه کسی هزینه پرداخت؟
\item \textbf{ساختار}: شرایط داخلی و بین‌المللی چگونه بود؟
\item \textbf{مکانیسم}: مداخله از چه نوعی بود (ارگانیک، ابزاری،
امپریال) و از چه مسیری (نرم تا سخت) گذشت؟
\item \textbf{نتایج}: ماتریس ۱۸ خانه‌ای چه تصویری ترسیم
می‌کند؟ نتیجه‌ی ژئوپولیتیک چه بود؟
\item \textbf{بازخورد}: چرخه‌ی روان‌شناختی چگونه عمل کرد؟
آیا چرخه شکسته شد یا بازتولید شد؟
\end{enumerate}
\end{tcolorbox}